\chapter{Dari Estimasi ke Difusi}
\label{ch:mlat}

Kuliah machine learning lanjut saya ikuti pada musim dingin 2025/26, di bawah Johannes
Brandstetter, yang risetnya sendiri berada di persimpangan deep learning dan simulasi fisika.
Delapan unitnya membentuk satu alur panjang: dari teori estimasi dan statistik Bayesian, lewat
VAE, persamaan diferensial biasa, parsial, dan stokastik, sampai model difusi, flow matching, dan
neural field. Tugas-tugasnya meliputi probabilistic PCA, score-based diffusion dan neural ODE, serta
variational diffusion model. Di antara semua kuliah, inilah yang paling langsung mengarah ke AI
untuk simulasi.

Benang merah bab ini sederhana untuk diucapkan, meski butuh beberapa halaman untuk dipahami:
model generatif modern adalah persamaan diferensial. Model difusi adalah persamaan difusi yang
dijalankan mundur. Flow matching adalah persamaan transport. Neural field adalah fungsi kontinu
yang bisa diatur oleh persamaan fisika. Untuk sampai ke sana, kita mulai dari pertanyaan yang
paling dasar: seberapa baik kita bisa menaksir sesuatu dari data.

\section{Teori estimasi}

Misalkan data $\mathcal D$ dibangkitkan oleh distribusi dengan parameter $\theta$ yang tidak
diketahui, dan kita membuat sebuah estimator $\hat\theta(\mathcal D)$. Karena data acak, estimator
pun acak. Kualitasnya diukur dengan mean squared error yang dirata-ratakan atas semua kemungkinan
data, dan MSE itu terurai menjadi varians dan bias kuadrat,
\begin{equation}
\mathrm{MSE}(\hat\theta)=\E\big[(\hat\theta-\theta)^2\big]=\mathrm{Var}(\hat\theta)+\big(\E[\hat\theta]-\theta\big)^2.
\end{equation}
Bentuknya sama dengan dekomposisi bias--varians di \cref{ch:data}, kali ini untuk parameter,
bukan untuk prediksi. Rata-rata sampel adalah estimator tak bias untuk rata-rata distribusi.
Varians sampel yang dibagi $N$ bias ke bawah, dan pembagian dengan $N-1$ menghilangkan bias itu.
Estimator tak bias dengan varians terkecil di antara semua estimator tak bias disebut
\istilah{minimum variance unbiased} (MVU).

Ada batas bawah mutlak untuk varians estimator tak bias mana pun. \istilah{Fisher information}
\begin{equation}
I(\theta)=-\E\Big[\frac{\partial^2}{\partial\theta^2}\log p(\mathcal D\mid\theta)\Big]
\end{equation}
mengukur seberapa tajam puncak log-likelihood di sekitar nilai yang benar. Puncak yang tajam
berarti data membawa banyak informasi tentang $\theta$. Batas Cramér--Rao menyatakan
$\mathrm{Var}(\hat\theta)\ge1/I(\theta)$ untuk setiap estimator tak bias. Untuk rata-rata distribusi
normal dengan varians $\sigma^2$ yang diketahui dan $N$ sampel, $I=N/\sigma^2$, sehingga varians
terkecil yang mungkin adalah $\sigma^2/N$. Rata-rata sampel mencapai batas ini, jadi ia estimator
yang efisien. Maximum likelihood, secara umum, tidak selalu tak bias untuk sampel kecil, tetapi
untuk sampel besar ia konsisten dan mencapai batas Cramér--Rao secara asimtotik.

Batas Cramér--Rao punya terjemahan langsung di dunia simulasi. Kalau kita ingin menaksir sebuah
parameter fisika, misalnya koefisien perpindahan panas, dari pengukuran temperatur, Fisher
information mengatakan apakah percobaan kita cukup informatif. Kalau log-likelihood hampir datar
terhadap parameter itu, tidak ada metode pintar yang bisa menaksirnya dengan baik. Masalahnya ada
pada desain percobaan, bukan pada algoritmenya.

\section{Cara pandang Bayesian}

Pendekatan frekuentis menganggap $\theta$ tetap dan datanya acak. Pendekatan Bayesian menganggap
$\theta$ sendiri tidak pasti dan mewakilinya dengan distribusi. Teorema Bayes memperbarui prior
menjadi posterior setelah melihat data. Untuk beberapa pasangan likelihood dan prior, posterior
punya bentuk yang sama dengan prior, dan pasangan seperti ini disebut \istilah{conjugate}. Contoh yang
dibahas kuliah adalah Gamma dan Poisson. Kalau jumlah kejadian per hari mengikuti Poisson dengan
laju $\lambda$, dan prior $\lambda\sim\mathrm{Gamma}(\alpha,\beta)$, maka setelah mengamati
$x_1,\dots,x_N$, posteriornya $\mathrm{Gamma}(\alpha+\sum_ix_i,\ \beta+N)$. Dengan prior
$\mathrm{Gamma}(2,1)$ dan pengamatan 3, 5, dan 4 kejadian, posteriornya $\mathrm{Gamma}(14,4)$ dengan
rata-rata $3{,}5$. Prior bertindak seperti dua kejadian yang ``sudah teramati'' dalam satu hari
semu, dan bobotnya memudar seiring datangnya data.

Estimasi MAP mengambil puncak posterior, dan dengan prior Gauss pada bobot ia menghasilkan weight
decay, seperti di \cref{ch:data}. Pendekatan Bayesian penuh tidak berhenti di puncak. Ia
memprediksi dengan merata-ratakan atas seluruh posterior, dan normalisasi posterior,
$p(\mathcal D)=\int p(\mathcal D\mid\theta)p(\theta)\dd\theta$, disebut \istilah{evidence}. Evidence
bisa dipakai untuk membandingkan model: model yang terlalu sederhana tidak bisa menjelaskan data,
sedangkan model yang terlalu fleksibel menyebar peluangnya ke terlalu banyak kemungkinan data,
sehingga keduanya mendapat evidence yang rendah.

\section{Variabel laten dan probabilistic PCA}

Banyak data lebih mudah dijelaskan kalau kita membayangkan ada sebab tersembunyi di baliknya.
\istilah{Latent variable model} menulis $p(\vx)=\int p(\vx\mid\vz)p(\vz)\dd\vz$, dengan $\vz$ yang
tidak pernah diamati. Contoh paling sederhana adalah probabilistic PCA \citep{tipping1999}:
\begin{equation}
\vx=\mW\vz+\bm\mu+\bm\varepsilon,\qquad \vz\sim\N(\bm0,\bm I_q),\qquad \bm\varepsilon\sim\N(\bm0,\sigma^2\bm I_D).
\end{equation}
Karena semuanya Gauss, distribusi marginalnya juga Gauss, $\vx\sim\N(\bm\mu,\mW\mW^\top+\sigma^2\bm I)$,
dan maximum likelihood punya solusi tertutup. Kolom $\mW$ terbentang oleh $q$ vektor eigen
utama dari kovarians sampel, diskalakan oleh $(\lambda_i-\sigma^2)^{1/2}$, dan $\sigma^2$ adalah
rata-rata nilai eigen yang dibuang. Ketika $\sigma^2\to0$, PCA biasa muncul kembali. Keuntungan
versi probabilistik adalah adanya likelihood, sehingga data yang hilang bisa ditangani, model
bisa dibandingkan, dan beberapa model bisa digabung.

Kalau $p(\vx\mid\vz)$ tidak lagi linear, misalnya berupa jaringan saraf, integral marginal tidak
bisa dihitung lagi, dan posterior $p(\vz\mid\vx)$ pun tidak. Dari sinilah VAE lahir.

\section{ELBO, sekali lagi dan lebih teliti}

\Cref{ch:arsitektur} memperkenalkan VAE secara intuitif. Kuliah ini menurunkannya. Ambil sembarang
distribusi $q(\vz\mid\vx)$. Log-likelihood bisa ditulis sebagai
\begin{equation}
\log p(\vx)=\underbrace{\E_{q(\vz\mid\vx)}\Big[\log\frac{p(\vx,\vz)}{q(\vz\mid\vx)}\Big]}_{\mathrm{ELBO}}
+\KL\big(q(\vz\mid\vx)\,\|\,p(\vz\mid\vx)\big).
\end{equation}
Identitas ini didapat dengan satu trik. Karena $\log p(\vx)$ tidak bergantung pada $\vz$, ia sama dengan
ekspektasinya di bawah $q$. Lalu tulis $p(\vx)=p(\vx,\vz)/p(\vz\mid\vx)$, kalikan dan bagi dengan $q(\vz\mid\vx)$
di dalam logaritma, dan pisahkan menjadi dua suku:
\[
\log p(\vx)=\E_q\Big[\log\frac{p(\vx,\vz)}{q(\vz\mid\vx)}\Big]+\E_q\Big[\log\frac{q(\vz\mid\vx)}{p(\vz\mid\vx)}\Big].
\]
Suku kedua adalah definisi divergensi KL. Divergensi KL tidak pernah negatif, sehingga suku pertama
adalah batas bawah untuk $\log p(\vx)$.
Memaksimalkan ELBO terhadap $q$ sama dengan mendekatkan $q$ ke posterior yang sebenarnya, dan
memaksimalkannya terhadap parameter model menaikkan likelihood. Kalau suku pertama diuraikan, kita
mendapat dua bagian yang sudah dikenal: rekonstruksi, $\E_q[\log p(\vx\mid\vz)]$, dan pencocokan
prior, $-\KL(q(\vz\mid\vx)\,\|\,p(\vz))$.

Divergensi KL sendiri layak dipahami. $\KL(q\|p)=\E_q[\log q-\log p]$ mengukur seberapa banyak
informasi yang hilang kalau $p$ dipakai untuk mewakili $q$. Ia tidak simetris. Meminimalkan
$\KL(q\|p)$ terhadap $q$ membuat $q$ menghindari daerah tempat $p$ kecil, sehingga $q$ cenderung
memilih satu puncak $p$ dan mengabaikan puncak lain.

Model campuran Gauss memberi contoh yang menjembatani. Data dianggap berasal dari salah satu dari
$K$ distribusi Gauss, dan label komponennya adalah variabel laten diskret. Algoritme EM untuk
model ini, yang sudah kita temui di \cref{ch:prob}, ternyata adalah maksimisasi ELBO bergantian:
langkah E memilih $q$ sama dengan posterior yang tepat, dan langkah M memperbarui parameter.
VAE melakukan hal yang sama, tetapi karena posteriornya tidak bisa dihitung, $q$ diwakili oleh
jaringan encoder, dan gradien dibawa melewati sampling dengan reparameterization trick.

Langkah terakhir menuju model difusi adalah VAE hierarkis. Alih-alih satu lapisan laten,
kita punya rantai $\vx\to\vz_1\to\vz_2\to\cdots\to\vz_T$, dan pada versi Markov setiap lapisan hanya
bergantung pada lapisan di sebelahnya. ELBO-nya terurai menjadi jumlah suku untuk setiap transisi.
Model difusi adalah VAE hierarkis Markov dengan dua keputusan istimewa: dimensi setiap laten sama
dengan dimensi data, dan encoder-nya tidak dipelajari, melainkan ditetapkan sebagai penambahan noise
Gauss sedikit demi sedikit \citep{luo2022}.

\section{Persamaan diferensial biasa, parsial, dan stokastik}

Sebelum masuk ke model difusi, kuliah berhenti sejenak untuk mengulang matematika yang akan dipakai.

\subsection{ODE}

Persamaan diferensial biasa, $\dot{\vx}=f(\vx,t)$ dengan $\vx(0)=\vx_0$, menggambarkan sistem yang
perubahannya ditentukan oleh keadaannya sekarang. Untuk sistem linear $\dot{\vx}=\mA\vx$, solusinya
$\vx(t)=e^{\mA t}\vx_0$, dan nilai eigen $\mA$ menentukan nasibnya: bagian real negatif berarti
peluruhan, bagian real positif berarti pertumbuhan, bagian imajiner berarti osilasi. Teorema
Picard--Lindelöf menjamin solusi ada dan tunggal kalau $f$ Lipschitz.

Kebanyakan ODE tidak punya solusi tertutup, jadi kita menyelesaikannya secara numerik. Metode Euler,
$\vx_{n+1}=\vx_n+h\,f(\vx_n,t_n)$, punya galat global sebanding dengan langkah $h$. Runge--Kutta orde
empat mengevaluasi $f$ empat kali per langkah dan galatnya sebanding dengan $h^4$. Untuk sistem kaku
(\istilah{stiff}), yaitu yang punya skala waktu cepat sekali dan lambat sekaligus, metode
eksplisit terpaksa memakai langkah kecil sekali agar stabil, dan metode implisit atau langkah
adaptif menjadi pilihan yang lebih baik. Topik ini dibahas lebih dalam di \cref{ch:fiskom}.

\subsection{PDE dan persamaan panas}

Persamaan diferensial parsial melibatkan turunan terhadap lebih dari satu variabel. Contoh yang
dibahas panjang lebar di kuliah adalah persamaan panas pada selang $(0,1)$,
\begin{equation}
\partial_tu=\kappa\,\partial_{xx}u,\qquad u(0,t)=u(1,t)=0.
\end{equation}
Pemisahan variabel memberi solusi
\begin{equation}
u(x,t)=\sum_{n=1}^\infty b_n\,e^{-\kappa(n\pi)^2t}\sin(n\pi x),
\end{equation}
dengan koefisien $b_n$ diambil dari kondisi awal. Setiap mode Fourier meluruh dengan laju yang
sebanding dengan kuadrat frekuensinya. Detail halus hilang paling cepat, dan setelah cukup lama
yang tersisa hanya gelombang paling panjang. Pada garis bilangan penuh, transformasi Fourier
memberi gambaran yang sama, $\hat u(k,t)=\hat u(k,0)\,e^{-\kappa k^2t}$, yang di ruang posisi berarti
konvolusi dengan kernel Gauss yang semakin lebar. Pengamatan ini akan berguna sebentar
lagi: menambahkan noise Gauss pada data sama dengan menjalankan persamaan panas pada distribusinya.

\subsection{Gerak Brown dan SDE}

Gerak Brown $W_t$ adalah batas acak dari jalan acak dengan langkah yang semakin kecil. Kenaikannya
independen dan berdistribusi normal, $W_{t+\Delta t}-W_t\sim\N(0,\Delta t)$. Simpangannya tumbuh
seperti $\sqrt{\Delta t}$, bukan $\Delta t$, sehingga lintasannya kontinu tetapi tidak bisa
diturunkan di mana pun. Sebuah jalan acak memberi gambaran: kalau setiap langkah waktu $\Delta t$ kita
melangkah ke kiri atau ke kanan sejauh $\sqrt{\Delta t}$, setelah waktu $t=n\Delta t$ variansnya
$n\Delta t=t$, tidak bergantung pada seberapa halus langkahnya. Kalau langkahnya sebanding $\Delta t$,
variansnya akan menyusut ke nol ketika langkah diperhalus. \istilah{Stochastic differential equation}
menambahkan suku acak ini ke ODE,
\begin{equation}
\dd\vx=f(\vx,t)\dd t+g(t)\dd W_t,
\end{equation}
dengan $f$ sebagai \istilah{drift} dan $g$ sebagai \istilah{diffusion}. Kalkulus untuk persamaan ini
punya satu kejutan. Karena $(\dd W)^2$ berukuran $\dd t$, bukan nol, aturan rantai mendapat suku
tambahan. Lemma Itô untuk fungsi skalar $\varphi(x)$ berbunyi
\begin{equation}
\dd\varphi=\Big(\varphi'f+\tfrac12g^2\varphi''\Big)\dd t+\varphi'g\dd W.
\end{equation}
Asal suku tambahan itu terlihat dari Taylor orde dua:
$\dd\varphi=\varphi'\dd x+\tfrac12\varphi''(\dd x)^2+\dots$. Dengan $\dd x=f\dd t+g\dd W$, kuadratnya memuat
$g^2(\dd W)^2$, dan karena $(\dd W)^2$ berukuran $\dd t$, suku itu bertahan, sedangkan $(\dd t)^2$ dan
$\dd t\,\dd W$ lebih kecil dan hilang. Contoh klasiknya gerak Brown geometris, $\dd S=\mu S\dd t+\sigma S\dd W$, model sederhana harga
saham. Dengan $\varphi=\log S$, lemma Itô memberi $\dd\log S=(\mu-\tfrac12\sigma^2)\dd t+\sigma\dd W$,
sehingga $S_t=S_0\exp\big((\mu-\tfrac12\sigma^2)t+\sigma W_t\big)$. Suku $-\tfrac12\sigma^2$ adalah jejak
kalkulus Itô yang tidak akan muncul dalam kalkulus biasa.

Distribusi $p(\vx,t)$ dari solusi SDE berevolusi menurut persamaan Fokker--Planck,
\begin{equation}
\partial_tp=-\nabla\cdot(f\,p)+\tfrac12g^2\,\Delta p.
\end{equation}
Suku pertama adalah transport oleh drift, suku kedua adalah difusi. Untuk $f=0$, ini persis
persamaan panas. Sebuah SDE yang bekerja pada partikel-partikel menjadi PDE yang bekerja pada
kerapatannya.

\section{Neural ODE}

Blok residual, $\vh_{l+1}=\vh_l+F(\vh_l)$, terlihat seperti satu langkah Euler. Kalau jumlah
lapisan dibuat tak hingga dan setiap langkahnya kecil, kita mendapat
\begin{equation}
\frac{\dd\vh}{\dd t}=f_\theta(\vh,t),
\end{equation}
sebuah \istilah{neural ODE} \citep{chen2018}. Lintasan maju dihitung dengan solver ODE apa pun,
termasuk yang langkahnya adaptif. Untuk gradien, kita bisa saja melakukan backprop melewati setiap
langkah solver, tetapi memorinya sebanding dengan jumlah langkah. Metode adjoint, yang dikenal lama
dalam teori kendali optimal, menghindarinya. Definisikan $\bm a(t)=\partial\Loss/\partial\vh(t)$.
Vektor ini memenuhi ODE-nya sendiri yang berjalan mundur,
\begin{equation}
\frac{\dd\bm a}{\dd t}=-\bm a^\top\frac{\partial f_\theta}{\partial\vh},\qquad
\frac{\dd\Loss}{\dd\theta}=-\int_{t_1}^{t_0}\bm a^\top\frac{\partial f_\theta}{\partial\theta}\dd t.
\end{equation}
Keadaan $\vh$, adjoint $\bm a$, dan gradien terhadap parameter diintegrasikan mundur bersama dalam
satu sistem yang diperbesar. Memorinya konstan terhadap jumlah langkah. Persamaan adjoint ini tidak
lain adalah delta backprop dari \cref{ch:dlbasic} dalam waktu kontinu.

Neural ODE cocok sekali untuk sistem fisika yang memang berevolusi secara kontinu. Data yang
diambil pada waktu tidak teratur bisa langsung dipakai, dan model dinamika yang dipelajari bisa
dijalankan dengan solver standar. Hubungan dengan model difusi juga erat: sampling deterministik
dari model difusi adalah menyelesaikan sebuah ODE.

\section{Model difusi}

Model difusi denoising \citep{sohldickstein2015,ho2020} menetapkan proses maju yang merusak data
selangkah demi selangkah,
\begin{equation}
q(\vx_t\mid\vx_{t-1})=\N\big(\sqrt{1-\beta_t}\,\vx_{t-1},\ \beta_t\bm I\big),
\end{equation}
dengan jadwal noise $\beta_t$ yang kecil. Karena komposisi Gauss tetap Gauss, kita bisa melompat
langsung dari data ke langkah mana pun: dengan $\bar\alpha_t=\prod_{s\le t}(1-\beta_s)$,
\begin{equation}
\vx_t=\sqrt{\bar\alpha_t}\,\vx_0+\sqrt{1-\bar\alpha_t}\,\bm\epsilon,\qquad \bm\epsilon\sim\N(\bm0,\bm I).
\end{equation}
Bentuk tertutup ini bisa diperiksa dengan dua langkah. Tulis $\alpha_t=1-\beta_t$. Satu langkah memberi
$\vx_1=\sqrt{\alpha_1}\vx_0+\sqrt{1-\alpha_1}\bm\epsilon_1$, dan langkah berikutnya
$\vx_2=\sqrt{\alpha_2}\vx_1+\sqrt{1-\alpha_2}\bm\epsilon_2=\sqrt{\alpha_1\alpha_2}\vx_0+\sqrt{\alpha_2(1-\alpha_1)}\bm\epsilon_1+\sqrt{1-\alpha_2}\bm\epsilon_2$.
Dua noise Gauss independen yang dijumlahkan tetap Gauss, dengan varians yang dijumlahkan:
$\alpha_2(1-\alpha_1)+(1-\alpha_2)=1-\alpha_1\alpha_2$. Jadi
$\vx_2=\sqrt{\alpha_1\alpha_2}\,\vx_0+\sqrt{1-\alpha_1\alpha_2}\,\bm\epsilon$, dan induksi memberi rumus umumnya.
Dengan jadwal $\beta_t$ yang naik linear dari $10^{-4}$ ke $0{,}02$ dalam seribu langkah, seperti pada
\citet{ho2020}, nilai $\bar\alpha_T$ sekitar $4\times10^{-5}$: sinyal aslinya tinggal kurang dari satu
persen dalam amplitudo. Setelah cukup banyak langkah, $\vx_T$ hampir murni noise. Proses mundur dipelajari oleh sebuah
jaringan. Dari ELBO VAE hierarkis, setelah beberapa langkah aljabar, loss pelatihannya menjadi
sederhana: tebak noise yang ditambahkan,
\begin{equation}
\Loss=\E_{t,\vx_0,\bm\epsilon}\big\|\bm\epsilon-\bm\epsilon_\theta(\vx_t,t)\big\|^2.
\end{equation}
Pada saat sampling, kita mulai dari noise murni dan berulang kali mengurangi noise yang ditebak
jaringan, ditambah sedikit noise baru, sampai kembali ke data.

Bayangkan setetes tinta jatuh ke segelas air. Tinta menyebar perlahan sampai airnya berwarna rata,
dan bentuk awal tetes itu tampak hilang. Model difusi belajar memutar video itu mundur: dari air
yang keruh merata, langkah demi langkah, sampai tetes tinta terbentuk kembali. Setiap langkah mundur
hanya perlu sedikit menebak, karena setiap langkah maju hanya sedikit merusak.

Kelemahan DDPM adalah lambat, karena sampling membutuhkan ratusan sampai ribuan langkah jaringan.
DDIM \citep{song2021ddim} menunjukkan bahwa jaringan yang sama bisa dipakai dengan proses sampling
deterministik yang melompati banyak langkah, tanpa pelatihan ulang. Variational diffusion models
\citep{kingma2021} merumuskan jadwal noise sebagai rasio sinyal terhadap noise yang bisa dipelajari dan
menunjukkan bahwa ELBO model difusi hanya bergantung pada jadwal itu melalui nilai ujung-ujungnya.

\section{Score dan persamaan diferensial stokastik}

\citet{song2021} menyatukan model difusi dalam bahasa SDE. Proses maju adalah SDE
$\dd\vx=f(\vx,t)\dd t+g(t)\dd W$ yang mengubah distribusi data menjadi Gauss. Hasil lama dari
\citet{anderson1982} menyatakan bahwa proses ini punya kebalikan dalam waktu, yang juga sebuah SDE,
\begin{equation}
\dd\vx=\big[f(\vx,t)-g(t)^2\,\nabla_{\vx}\log p_t(\vx)\big]\dd t+g(t)\dd\bar W,
\end{equation}
yang berjalan dari $t=T$ ke $t=0$. Satu-satunya besaran yang tidak diketahui adalah \istilah{score}
$\nabla_{\vx}\log p_t(\vx)$, yaitu arah tempat kerapatan data naik paling cepat. Jaringan
$s_\theta(\vx,t)$ dilatih untuk menebak score ini. Melatihnya ternyata mudah berkat
\istilah{denoising score matching} \citep{vincent2011}: kalau $\tilde\vx=\vx+\sigma\bm\epsilon$, score
dari distribusi bersyarat $p(\tilde\vx\mid\vx)$ adalah $-(\tilde\vx-\vx)/\sigma^2=-\bm\epsilon/\sigma$,
sehingga menebak score sama dengan menebak noise, persis seperti DDPM.

Ada juga ODE yang menghasilkan distribusi marginal yang sama di setiap waktu, yaitu
\istilah{probability flow ODE},
\begin{equation}
\dd\vx=\big[f(\vx,t)-\tfrac12g(t)^2\nabla_{\vx}\log p_t(\vx)\big]\dd t.
\end{equation}
Mengapa SDE dan ODE bisa menghasilkan marginal yang sama? Persamaan Fokker--Planck menjawabnya.
Suku difusi $\tfrac12g^2\Delta p$ bisa ditulis sebagai $\nabla\cdot(\tfrac12g^2\,p\,\nabla\log p)$,
sehingga difusi pada kerapatan bisa diganti dengan transport deterministik dengan kecepatan
$-\tfrac12g^2\nabla\log p$. Partikel-partikel individu bergerak berbeda, tetapi kerumunannya
berevolusi dengan cara yang sama. Kuliah merangkum bahwa ketiga cara melatih model difusi, yaitu
menebak data bersih, menebak noise, dan menebak score, setara sampai faktor skala. Untuk
membangkitkan sampel bersyarat, misalnya gambar dari teks, score ditambah dengan gradien sebuah
penilai syarat, atau syarat langsung diberikan sebagai input ke jaringan.

\section{Flow matching}

Kalau ujung-ujungnya kita menyelesaikan sebuah ODE, mengapa tidak langsung belajar ODE-nya?
Flow matching \citep{lipman2023} melakukan hal itu. Kita mencari medan kecepatan $v_\theta(\vx,t)$
sehingga ODE $\dot{\vx}=v_\theta(\vx,t)$ membawa sampel dari distribusi sederhana $p_0$, misalnya
Gauss, ke distribusi data $p_1$. Pemetaan yang dihasilkan disebut \istilah{flow} dalam arti
matematis, dan barisan distribusi antara $p_t$ disebut \istilah{probability path}.

Kesulitannya, medan kecepatan yang benar untuk seluruh distribusi tidak diketahui. Kunci flow
matching adalah bekerja secara bersyarat. Untuk satu pasangan titik, sampel noise $\vx_0$ dan sampel
data $\vx_1$, kita bisa memilih jalur yang paling sederhana, garis lurus
$\vx_t=(1-t)\vx_0+t\vx_1$, yang kecepatannya konstan, $\vx_1-\vx_0$. Loss \istilah{conditional flow
matching} memaksa jaringan meniru kecepatan bersyarat itu,
\begin{equation}
\Loss_{\mathrm{CFM}}=\E_{t,\vx_0,\vx_1}\big\|v_\theta(\vx_t,t)-(\vx_1-\vx_0)\big\|^2,
\end{equation}
dan hasil utamanya adalah bahwa gradien loss ini sama dengan gradien loss untuk medan kecepatan
marginal yang sebenarnya. Kita tidak pernah perlu mengetahui medan marginal itu. Gagasan jalur lurus
juga dikembangkan secara independen sebagai \istilah{rectified flow} \citep{liu2023}.

Analogi yang saya pakai adalah rombongan penumpang yang harus dipindahkan dari terminal ke berbagai
alamat di kota. Kita tidak tahu jalur terbaik untuk seluruh rombongan, tetapi kita tahu jalur lurus
untuk setiap pasangan terminal dan alamat. Kalau sopir belajar rata-rata dari jutaan perjalanan
lurus seperti itu, ia akhirnya tahu arah yang tepat di setiap persimpangan untuk seluruh arus
penumpang. Persamaan kontinuitas yang mengatur probability path, $\partial_tp_t+\nabla\cdot(p_tv_t)=0$,
adalah persamaan kekekalan massa yang sama dengan di mekanika fluida.

\section{Neural field}

Unit terakhir kuliah meninggalkan model generatif dan membahas cara lain merepresentasikan data.
Dalam fisika, medan adalah besaran yang punya nilai di setiap titik ruang dan waktu: temperatur,
tekanan, kecepatan. \istilah{Neural field} adalah jaringan saraf yang mewakili medan seperti itu
secara langsung, $f_\theta:(\vx,t)\mapsto\text{nilai}$ \citep{xie2022}. Alih-alih menyimpan nilai di
grid, kita menyimpan bobot jaringan. Representasinya kontinu, tidak terikat resolusi, dan bisa
diturunkan di mana saja dengan autodiff.

Bentuk tiga dimensi bisa diwakili sebagai occupancy field, yaitu peluang sebuah titik berada di
dalam benda, atau sebagai \istilah{signed distance field} $f(\vx)$, jarak bertanda dari titik ke
permukaan, positif di luar dan negatif di dalam \citep{park2019}. Permukaannya adalah himpunan nol
$f(\vx)=0$, dan bisa diekstrak dengan algoritme seperti marching cubes. Fungsi jarak yang sejati
memenuhi persamaan Eikonal, $\|\nabla f(\vx)\|=1$. Menambahkan $(\|\nabla f_\theta\|-1)^2$ ke loss di
titik-titik acak memaksa jaringan menjadi fungsi jarak yang sah \citep{gropp2020}. Ini adalah loss
yang diinformasikan fisika: sebuah persamaan diferensial parsial ditegakkan di dalam loss.

NeRF, yang sudah kita lihat di \cref{ch:cv}, adalah neural field untuk kerapatan dan warna yang
dilatih melalui persamaan volume rendering. Kuliah menutup dengan pernyataan yang menjembatani ke
Bagian III: dari sudut pandang ini, PINN hanyalah neural field yang diatur oleh persamaan fisika.
Bedanya dengan NeRF hanya pada persamaan apa yang ditegakkan.

Satu masalah praktis menghubungkan semuanya. MLP biasa dengan aktivasi ReLU cenderung belajar
komponen frekuensi rendah lebih dulu dan kesulitan mewakili detail halus, gejala yang disebut
\istilah{spectral bias} \citep{rahaman2019}. Dua obat yang populer adalah memetakan input dengan
fungsi sinus dan kosinus berbagai frekuensi sebelum masuk jaringan \citep{tancik2020}, atau memakai
fungsi sinus sebagai aktivasi seperti pada SIREN \citep{sitzmann2020}. Spectral bias juga menjadi
salah satu alasan PINN sulit mempelajari solusi yang berosilasi cepat (\cref{ch:pinn}).

\section{Di luar kelas}

Alat-alat dari bab ini menyebar jauh. Model difusi membangkitkan gambar dan video (\cref{ch:terbaru}), dan
modul struktur AlphaFold3 adalah model difusi yang membangkitkan koordinat atom (\cref{ch:protein}). Teori
estimasi dan batas Cramér--Rao adalah bahasa sehari-hari pemrosesan sinyal (\cref{ch:sinyal}).
Probabilistic PCA dan analisis faktor dipakai untuk meringkas data ekspresi gen yang berisik
(\cref{ch:genomik}). Neural ODE cocok untuk data medis yang diukur pada waktu tidak teratur, seperti kunjungan
pasien yang jaraknya berbeda-beda.

Bagi Bagian V, bab ini memberi hampir semua alatnya. Neural ODE adalah cara alami untuk mempelajari
dinamika sistem fisika dari data. Neural field adalah kerangka PINN. Model difusi dan flow matching
dipakai untuk membangkitkan medan turbulen yang realistis, untuk memperbaiki prediksi jangka panjang
model pengganti \citep{lippe2023}, dan untuk prakiraan cuaca probabilistik \citep{price2025}. Bahkan
cara berpikirnya pun berasal dari fisika: model difusi adalah persamaan panas pada ruang
distribusi, dan flow matching adalah persamaan kontinuitas. Batas Cramér--Rao mengingatkan bahwa
ketika kita mencoba menyimpulkan parameter fisika dari data, informasi yang terkandung di dalam data
membatasi apa yang bisa dicapai oleh metode mana pun.

\sumberbab{Bab ini mengikuti delapan unit kuliah Machine Learning: Advanced Techniques: teori estimasi,
statistik Bayesian, probabilistic PCA, VAE dan ELBO, ODE, PDE, dan SDE, neural ODE, model difusi
(DDPM, DDIM), score-based model dan flow matching, Fokker--Planck, dan neural field. Bacaan utamanya
adalah makalah-makalah aslinya: \citet{tipping1999}, \citet{kingma2014}, \citet{chen2018},
\citet{ho2020}, \citet{song2021}, \citet{lipman2023}, dan ulasan \citet{luo2022}.}
